\section*{Hoofdstuk \ref{ch:b2:huckel} --- Hückeltheorie en geconjugeerde systemen}

\begin{solution}{exo:b2:huckel:1}
Niveaus van allyl ($n = 3$): $\alpha + 1.414\beta$, $\alpha$, $\alpha - 1.414\beta$. Kation (2 elektronen):
$2\alpha + 2.828\beta$; radicaal (3): $3\alpha + 2.828\beta$; anion (4): $4\alpha + 2.828\beta$. De extra
elektronen gaan naar het niet-bindende niveau en veranderen alleen het $\alpha$-deel.
\end{solution}

\begin{solution}{exo:b2:huckel:2}
$2(\alpha + 1.618\beta) + 2(\alpha + 0.618\beta) = 4\alpha + 4.472\beta$.
\end{solution}

\begin{solution}{exo:b2:huckel:3}
Benzeen (6), cyclopentadienylanion (6), tropylium (6): \omterm{def:b2:huckel:aromatic}{aromatisch}. Cyclopentadienylkation
(4), cyclobutadieen (4), vlak cyclo-octatetraeen (8): \omterm{def:b2:huckel:aromatic}{antiaromatisch} als het vlak is
(echt cyclo-octatetraeen is kuipvormig en \omterm{def:b2:huckel:aromatic}{niet-aromatisch}).
\end{solution}

\begin{solution}{exo:b2:huckel:4}
Een achthoek in de cirkel: niveaus $\alpha + 2\beta$, $\alpha + 1.414\beta$ (tweemaal), $\alpha$ (tweemaal),
$\alpha - 1.414\beta$ (tweemaal), $\alpha - 2\beta$. Acht elektronen: twee in het laagste, vier in het
eerste paar, één in elke orbitaal van het niet-bindende paar.
\end{solution}

\begin{solution}{exo:b2:huckel:5}
$4.472\beta - 4\beta = 0.472\beta$, dus $0.472 \times 75 = \qty{35}{kJ/mol}$, tegenover ongeveer
\qty{10}{kJ/mol} uit de hydrogeneringsgegevens: het model van Hückel, met $\beta$ aangepast aan benzeen,
overschat de conjugatie van open ketens.
\end{solution}

\begin{solution}{exo:b2:huckel:6}
$p_{12} = p_{34} = 0.894$, $p_{23} = 0.447$: de eindbindingen C1--C2 en C3--C4 zijn de
kortste, de centrale binding is langer, maar korter dan een enkelvoudige binding.
\end{solution}

\begin{solution}{exo:b2:huckel:7}
Anion, zes elektronen: $2 \times 2 + 4 \times 0.618$, $E_\pi = 6\alpha + 6.472\beta$, een gesloten schil:
\omterm{def:b2:huckel:aromatic}{aromatisch}. Kation, vier: $2 \times 2 + 2 \times 0.618$, $4\alpha + 5.236\beta$, met één elektron
in elke orbitaal van het ontaarde paar: \omterm{def:b2:huckel:aromatic}{antiaromatisch}, zeer instabiel.
\end{solution}

\begin{solution}{exo:b2:huckel:8}
Zevenring: $\alpha + 2\beta$, $\alpha + 1.247\beta$ (tweemaal), daarna antibindende niveaus. Zes
elektronen, zoals in het kation \ce{C7H7+}, vullen precies de bindende niveaus: een \omterm{def:b2:huckel:aromatic}{aromatisch}
kation. De verbinding ioniseert gemakkelijk tot tropylium en bromide.
\end{solution}

\begin{solution}{exo:b2:huckel:9}
$1.802 + 1.247 + 0.445 - 0.445 - 1.247 - 1.802 = 0$: de som van de energieën is $6\alpha$.
\end{solution}

\begin{solution}{exo:b2:huckel:10}
$n = 6$: $m_k = 2\cos(k\pi/7) = 1.802$, $1.247$, $0.445$, $-0.445$, $-1.247$, $-1.802$. Zes
elektronen: $E_\pi = 6\alpha + 2(1.802 + 1.247 + 0.445)\beta = 6\alpha + 6.988\beta$;
delokalisatie $0.988\beta$.
\end{solution}

\begin{solution}{exo:b2:huckel:11}
In het vierkant zitten de twee elektronen met de hoogste energie elk in één orbitaal van een
ontaard niet-bindend paar. Het uitrekken van twee tegenover elkaar liggende bindingen heft de ontaarding op:
de ene orbitaal daalt, de andere stijgt, en beide elektronen paren in de laagste, wat
de energie verlaagt. Het molecuul neemt de vorm aan van een rechthoek met twee gelokaliseerde dubbele
bindingen; het blijft \omterm{def:b2:huckel:aromatic}{antiaromatisch} en uiterst reactief.
\end{solution}

\begin{solution}{exo:b2:huckel:12}
Matrix (in eenheden van $\beta$, met $\alpha$ als oorsprong) $\begin{pmatrix}0 & 1\\ 1 & 1\end{pmatrix}$:
$m = 1.618$ en $-0.618$, $E = \alpha + 1.618\beta$ (bindend) en $\alpha - 0.618\beta$. Voor de
\omterm{def:b2:diatomic-mos:bonding}{bindende orbitaal} is $c_X = 1.618c_C$, dus $c_C^2 = 0.276$, $c_X^2 = 0.724$; met twee elektronen
is $q_C = 0.55$, $q_X = 1.45$: de $\pi$-dichtheid wordt naar X getrokken.
\end{solution}

\begin{solution}{pb:b2:huckel:1}
\textbf{1.} $-123.1 - (-4.3) = \qty{-118.8}{kJ/mol}$.
\textbf{2.} $-123.1 - 104.6 = \qty{-227.7}{kJ/mol}$.
\textbf{3.} $-123.1 - 82.9 = \qty{-206.0}{kJ/mol}$.
\textbf{4.} $3 \times (-118.8) = \qty{-356.4}{kJ/mol}$.
\textbf{5.} $356.4 - 206.0 = \qty{150}{kJ/mol}$.
\textbf{6.} Het dieen geeft $2 \times 118.8 - 227.7 = \qty{10}{kJ/mol}$ minder vrij dan twee
geïsoleerde dubbele bindingen: conjugatie van een open keten levert een kleine winst op, de gesloten
ring van benzeen vijftien keer meer.
\textbf{7.} Een matrix $6 \times 6$ met 1 tussen elk atoom en zijn twee buren (1--2, 2--3,
\dots, 6--1) en elders 0.
\textbf{8.} $m = 2\cos(2\pi k/6)$: $2$, $1$, $1$, $-1$, $-1$, $-2$: $E = \alpha + 2\beta$, $\alpha + \beta$ (tweemaal),
$\alpha - \beta$ (tweemaal), $\alpha - 2\beta$.
\textbf{9.} Twee elektronen in $\alpha + 2\beta$, vier in het paar $\alpha + \beta$.
\textbf{10.} $E_\pi = 6\alpha + 8\beta$.
\textbf{11.} $3(2\alpha + 2\beta) = 6\alpha + 6\beta$.
\textbf{12.} $2\beta$, dus $2|\beta|$ stabilisatie.
\textbf{13.} $2 + 1 + 1 - 1 - 1 - 2 = 0$: de zes niveaus tellen op tot $6\alpha$.
\textbf{14.} $2|\beta| = \qty{150}{kJ/mol}$: $|\beta| = \qty{75}{kJ/mol}$.
\textbf{15.} $0.472 \times 75 = \qty{35}{kJ/mol}$, tegenover \qty{10}{kJ/mol}: het model dat
op benzeen is aangepast, overschat de open keten.
\textbf{16.} De bezette orbitalen kunnen genomen worden als $\psi_1 = (1, 1, 1, 1, 1, 1)/\sqrt6$,
$\psi_2 = (2, 1, -1, -2, -1, 1)/\sqrt{12}$ en $\psi_3 = (0, 1, 1, 0, -1, -1)/2$. Voor de binding
1--2: $p_{12} = 2(\frac16 + \frac{2}{12} + 0) = \frac23$; door symmetrie heeft elke binding $p = 2/3$.
\textbf{17.} Tussen de eindbindingen (0.894) en de centrale binding (0.447) van butadieen.
\textbf{18.} De zes bindingen zijn door symmetrie gelijkwaardig; hun lengte, \qty{139.7}{pm}, ligt
tussen de dubbele binding van etheen (\qty{133.9}{pm}) en de enkelvoudige binding van ethaan
(\qty{153.6}{pm}).
\textbf{19.} $0.988|\beta| = \qty{74}{kJ/mol}$, de helft van die van benzeen: het sluiten van de ring verdubbelt de
\omterm{def:b2:huckel:delocalisation}{delokalisatie-energie}.
\textbf{20.} $\alpha + 2\beta$, $\alpha + 1.414\beta$ (tweemaal), $\alpha$ (tweemaal), $\alpha - 1.414\beta$ (tweemaal),
$\alpha - 2\beta$.
\textbf{21.} De laatste twee elektronen gaan elk naar één orbitaal van het niet-bindende paar: een open
schil, \omterm{def:b2:huckel:aromatic}{antiaromatisch}.
\textbf{22.} Het plooit tot een kuip: de p-orbitalen van naburige dubbele bindingen zijn niet langer
parallel, de bindingen alterneren, en het molecuul gedraagt zich als een polyeen.
\textbf{23.} Vier elektronen: hetzelfde half gevulde paar, erger omdat het molecuul
vlakheid niet kan vermijden; het vervormt tot een rechthoek en is uiterst reactief.
\textbf{24.} Een vlak monocyclisch geconjugeerd systeem is \omterm{def:b2:huckel:aromatic}{aromatisch} met $4N + 2$ $\pi$-elektronen,
\omterm{def:b2:huckel:aromatic}{antiaromatisch} met $4N$.
\textbf{25.} $\boldsymbol{|\beta| \approx \qty{75}{kJ/mol}}$.
\end{solution}
